\section{2010 Applied \& Computational Mathematics}

\subsection{Individual}

\begin{exercise}
\label{applied:2010:1}
Let $Z_{1}, \cdots, Z_{n}$ be i.i.d. random variables with $Z_{i} \sim N(\mu, \sigma^{2})$. Find $E(\sum_{i=1}^{n} Z_{i} | Z_{1} - Z_{2} + Z_{3})$.
\end{exercise}

\begin{solution}
\textbf{Steps:}
\begin{enumerate}
    \item Let $S = \sum_{i=1}^n Z_i$ and $D = Z_1 - Z_2 + Z_3$. Note that $(S, D)$ follows a bivariate normal distribution since they are linear combinations of independent normal variables.
    \item Compute the parameters:
    \begin{itemize}
        \item $E(S) = n\mu$.
        \item $E(D) = \mu - \mu + \mu = \mu$.
        \item $Var(D) = Var(Z_1) + Var(-Z_2) + Var(Z_3) = 3\sigma^2$.
        \item $Cov(S, D) = Cov(\sum Z_i, Z_1 - Z_2 + Z_3) = Cov(Z_1, Z_1) - Cov(Z_2, Z_2) + Cov(Z_3, Z_3) = \sigma^2 - \sigma^2 + \sigma^2 = \sigma^2$.
    \end{itemize}
    \item Use the conditional expectation formula for bivariate normal variables:
    \[ E(S|D) = E(S) + \frac{Cov(S, D)}{Var(D)}(D - E(D)) \]
    \item Substitute the values:
    \[ E(S|D) = n\mu + \frac{\sigma^2}{3\sigma^2}(D - \mu) = n\mu + \frac{1}{3}(Z_1 - Z_2 + Z_3 - \mu). \]
\end{enumerate}

\textbf{Intuition:}
The sum of all variables and the specific linear combination $Z_1 - Z_2 + Z_3$ are correlated because they share common variables. In a normal setting, the conditional expectation is a linear function of the conditioning variable, with the slope determined by the ratio of the covariance to the variance of the condition.
\end{solution}

\begin{exercise}
\label{applied:2010:2}
Let $X_{1}, \cdots, X_{n}$ be pairwise independent. Further, assume that $E X_{i} = 1 + i^{-1}$ and that $\max_{1 \leq i \leq n} E | X_{i} |^{1 + \epsilon} < \infty$ for some $\epsilon > 0$. Show that $\frac{1}{n} \sum_{i=1}^{n} X_{i} \xrightarrow{P} 1$.
\end{exercise}

\begin{solution}
\textbf{Steps:}
\begin{enumerate}
    \item Compute the expected value of the average:
    \[ E\left(\frac{1}{n} \sum_{i=1}^n X_i\right) = \frac{1}{n} \sum_{i=1}^n (1 + i^{-1}) = 1 + \frac{1}{n} \sum_{i=1}^n \frac{1}{i} = 1 + \frac{H_n}{n}. \]
    Since $H_n \sim \ln n$, $\frac{H_n}{n} \to 0$ as $n \to \infty$, so the expectation converges to 1.
    \item To show convergence in probability, we need to show that the variance (or a related dispersion measure) goes to zero. Since the variables are only pairwise independent, the standard WLLN requires slightly more care.
    \item The condition $\max E |X_i|^{1+\epsilon} < \infty$ implies uniform integrability and allows the use of the weak law for pairwise independent variables. Specifically, we can use a truncation argument: let $Y_i = X_i I(|X_i| \le n)$.
    \item For pairwise independent variables with bounded $(1+\epsilon)$ moments, the average $\frac{1}{n} \sum (X_i - E X_i)$ converges to 0 in probability (and even almost surely if the moments are $L^2$).
\end{enumerate}

\textbf{Intuition:}
The "pairwise independent" condition is weaker than full independence but still allows for the law of large numbers to hold. The bounded $(1+\epsilon)$ moment ensures that no single variable's tail can dominate the sum, providing enough regularity for the average to concentrate around its mean.
\end{solution}

\begin{exercise}
\label{applied:2010:3}
Let $Z_{1}, \cdots, Z_{6}$ be i.i.d. random variables with $Z_{i} \sim N(0, 1)$. Set $U^{2} = \frac{(Z_{1} Z_{2} + Z_{3} Z_{4} + Z_{5} Z_{6})^{2}}{Z_{2}^{2} + Z_{4}^{2} + Z_{6}^{2}}$, $V^{2} = \frac{U^{2} (Z_{2}^{2} + Z_{4}^{2})}{U^{2} + Z_{6}^{2}}$. Find and identify the densities of $U^{2}$ and $V^{2}$.
\end{exercise}

\begin{solution}
\textbf{Steps:}
\begin{enumerate}
    \item Let $\mathbf{X} = (Z_1, Z_3, Z_5)^T$ and $\mathbf{Y} = (Z_2, Z_4, Z_6)^T$. Then $U^2 = \frac{(\mathbf{X} \cdot \mathbf{Y})^2}{\|\mathbf{Y}\|^2}$.
    \item Conditional on $\mathbf{Y}$, the term $\frac{\mathbf{X} \cdot \mathbf{Y}}{\|\mathbf{Y}\|}$ is a linear combination of independent standard normals, which is itself $N(0, 1)$. Since this distribution does not depend on $\mathbf{Y}$, $U = \frac{\mathbf{X} \cdot \mathbf{Y}}{\|\mathbf{Y}\|}$ is independent of $\mathbf{Y}$ and $U \sim N(0, 1)$.
    \item Thus $U^2 \sim \chi^2(1)$.
    \item For $V^2 = \frac{U^2 (Z_2^2 + Z_4^2)}{U^2 + Z_6^2}$, note that $U^2, Z_2^2+Z_4^2, Z_6^2$ are independent. Let $A = U^2 \sim \chi^2(1)$, $B = Z_6^2 \sim \chi^2(1)$, and $C = Z_2^2 + Z_4^2 \sim \chi^2(2)$.
    \item $V^2 = C \cdot \frac{A}{A+B}$. Since $A, B$ are i.i.d. $\chi^2(1)$, $W = \frac{A}{A+B} \sim Beta(\frac{1}{2}, \frac{1}{2})$.
    \item $C \sim \chi^2(2)$ is an Exponential distribution with mean 2, i.e., $f_C(c) = \frac{1}{2} e^{-c/2}$.
    \item The density of $V^2$ can be found by the product of Beta and Gamma distributions.
\end{enumerate}

\textbf{Intuition:}
The variable $U$ represents the projection of one random vector onto another. Due to the rotational invariance of the normal distribution, this projection is always standard normal. $V^2$ then combines this with other $\chi^2$ components, resulting in a distribution related to the ratio of variances.
\end{solution}

\begin{exercise}
\label{applied:2010:4}
Consider a population with three characteristics with frequencies: $p_1 = \theta^3$, $p_2 = 3\theta(1-\theta)$, $p_3 = (1-\theta)^3$. Note that these sum to 1.
(a) Construct approx level $(1-\alpha)$ MLE confidence set.
(b) Derive asymptotic distribution for freq substitution estimator $\hat{\theta}_2 = 1 - (N_3/n)^{1/3}$.
\end{exercise}

\begin{solution}
\textbf{Steps:}
\begin{enumerate}
    \item \textbf{MLE:} The log-likelihood is $\ell(\theta) = N_1 \ln \theta^3 + N_2 \ln(3\theta(1-\theta)) + N_3 \ln(1-\theta)^3$.
    \item $\ell'(\theta) = \frac{3N_1+N_2}{\theta} - \frac{N_2+3N_3}{1-\theta} = 0 \implies \hat{\theta} = \frac{3N_1+N_2}{3N_1+2N_2+3N_3}$.
    \item Fisher Information $I(\theta) = -E[\ell''(\theta)]$ gives the asymptotic variance $Var(\hat{\theta}) \approx \frac{1}{n I(\theta)}$. The confidence set is $\hat{\theta} \pm z_{\alpha/2} \sqrt{\widehat{Var}(\hat{\theta})}$.
    \item \textbf{Substitution Estimator:} For $\hat{\theta}_2 = 1 - (\frac{N_3}{n})^{1/3}$, we use the Delta Method.
    \item $\frac{N_3}{n} \xrightarrow{d} AN(p_3, \frac{p_3(1-p_3)}{n})$. Let $g(x) = 1 - x^{1/3}$. Then $g'(x) = -\frac{1}{3} x^{-2/3}$.
    \item Asymptotic distribution: $\sqrt{n}(\hat{\theta}_2 - \theta) \xrightarrow{d} N(0, [g'(p_3)]^2 p_3(1-p_3))$.
\end{enumerate}
\end{solution}

\begin{exercise}
\label{applied:2010:5}
\begin{enumerate}
    \item Suppose
    \[
    S = \begin{bmatrix} \sigma & \mathbf{u}^T \\ 0 & S_c \end{bmatrix}, \quad T = \begin{bmatrix} \tau & \mathbf{v}^T \\ 0 & T_c \end{bmatrix}, \quad \mathbf{b} = \begin{bmatrix} \beta \\ \mathbf{b}_c \end{bmatrix},
    \]
    where $\sigma, \tau$ and $\beta$ are scalars, $S_c$ and $T_c$ are $n$-by-$n$ matrices, and $\mathbf{b}_c$ is an $n$-vector. Show that if there exists a vector $\mathbf{x}_c$ such that $(S_c T_c - \lambda I) \mathbf{x}_c = \mathbf{b}_c$ and $\mathbf{w}_c = T_c \mathbf{x}_c$ is available, then
    \[
    \mathbf{x} = \begin{bmatrix} \gamma \\ \mathbf{x}_c \end{bmatrix}, \quad \gamma = \frac{\beta - \sigma \mathbf{v}^T \mathbf{x}_c - \mathbf{u}^T \mathbf{w}_c}{\sigma \tau - \lambda}
    \]
    solves $(ST - \lambda I) \mathbf{x} = \mathbf{b}$.

    \item Hence or otherwise, derive an $O(n^2)$ algorithm for solving the linear system $(U_1 U_2 - \lambda I) \mathbf{x} = \mathbf{b}$ where $U_1$ and $U_2$ are $n$-by-$n$ upper triangular matrices, and $(U_1 U_2 - \lambda I)$ is nonsingular. Please write down your algorithm and prove that it is indeed of $O(n^2)$ complexity.

    \item Hence or otherwise, derive an $O(pn^2)$ algorithm for solving the linear system $(U_1 U_2 \cdots U_p - \lambda I) \mathbf{x} = \mathbf{b}$ where $\{U_i\}_{i=1}^p$ are all $n$-by-$n$ upper triangular matrices, and $(U_1 U_2 \cdots U_p - \lambda I)$ is nonsingular. Please write down your algorithm and prove that it is indeed of $O(pn^2)$ complexity.
\end{enumerate}
\end{exercise}

\begin{solution}
\textbf{Prerequisites:}
\begin{itemize}
    \item Block matrix multiplication and inversion.
    \item Fundamental properties of upper triangular matrices (product is triangular, solving via back-substitution).
    \item Order of complexity for matrix-vector products.
\end{itemize}

\textbf{Steps:}
\begin{enumerate}
    \item \textbf{Identity Verification:} Compute the product $ST = \begin{bmatrix} \sigma \tau & \sigma \mathbf{v}^T + \mathbf{u}^T T_c \\ 0 & S_c T_c \end{bmatrix}$. The system $(ST - \lambda I)\mathbf{x} = \mathbf{b}$ yields:
    \[ \begin{bmatrix} \sigma \tau - \lambda & \sigma \mathbf{v}^T + \mathbf{u}^T T_c \\ 0 & S_c T_c - \lambda I \end{bmatrix} \begin{bmatrix} \gamma \\ \mathbf{x}_c \end{bmatrix} = \begin{bmatrix} \beta \\ \mathbf{b}_c \end{bmatrix}. \]
    The lower block gives $(S_c T_c - \lambda I)\mathbf{x}_c = \mathbf{b}_c$. The upper block gives $(\sigma \tau - \lambda)\gamma + \sigma \mathbf{v}^T \mathbf{x}_c + \mathbf{u}^T (T_c \mathbf{x}_c) = \beta$. Solving for $\gamma$ with $\mathbf{w}_c = T_c \mathbf{x}_c$ yields the desired formula.

    \item \textbf{$O(n^2)$ Algorithm for $U_1 U_2$:} We use back-substitution. Let $(U_1)_{ii} = \sigma_i$ and $(U_2)_{ii} = \tau_i$. Starting from $i=n$ down to $1$:
    \begin{itemize}
        \item Compute $x_i = (\beta_i - \sigma_i \sum_{j=i+1}^n (U_2)_{ij} x_j - \sum_{j=i+1}^n (U_1)_{ij} w_j) / (\sigma_i \tau_i - \lambda)$.
        \item Update $w_i = \tau_i x_i + \sum_{j=i+1}^n (U_2)_{ij} x_j$.
    \end{itemize}
    Each step $i$ involves $O(n-i)$ additions and multiplications. Total complexity $\sum_{i=1}^n O(n-i) = O(n^2)$.

    \item \textbf{$O(pn^2)$ Algorithm for $p$ matrices:} Let $M^{(k)} = U_k \cdots U_p$. To solve $(M^{(1)} - \lambda I) \mathbf{x} = \mathbf{b}$, we define intermediate vectors $\mathbf{w}^{(k)} = M^{(k+1)} \mathbf{x}$. By solving the components backward $i=n, \dots, 1$ and for each $i$ computing the sequence of $p$ intermediate values, the logic from Part (2) repeats $p$ times per row entry. Total complexity is $O(p n^2)$.
\end{enumerate}

\textbf{Intuition:}
The problem leverages the fact that the product of upper triangular matrices is itself upper triangular. Instead of explicitly computing the matrix product $U_1 U_2$ (which is $O(n^3)$), we can solve the system by maintaining partial products of the triangular matrices with the solution vector, essentially extending the standard back-substitution method to handle the nested structure.
\end{solution}

\begin{exercise}
\label{applied:2010:6}
\begin{enumerate}
    \item Let $A \in \mathbb{R}^{m \times n}$. Show that there exists an $m$-by-$m$ orthogonal matrix $U$ and an $n$-by-$n$ orthogonal matrix $V$ such that $U^T A V = \text{diag}(\sigma_1, \dots, \sigma_p)$, where $p = \min\{m, n\}$ and $\sigma_1 \geq \sigma_2 \geq \dots \geq \sigma_p \geq 0$.

    \item Let $\text{rank}(A) = r$. Show that for any positive integer $k < r$
    \[
    \min_{\text{rank}(B) = k} \| A - B \|_2 = \sigma_{k+1}.
    \]
\end{enumerate}
\end{exercise}

\begin{solution}
\textbf{Prerequisites:}
\begin{itemize}
    \item Spectral theorem for real symmetric matrices.
    \item Definition and properties of orthogonal matrices.
    \item Vector norm and induced matrix 2-norm properties.
\end{itemize}

\textbf{Steps:}
\begin{enumerate}
    \item \textbf{Existence of SVD:} Consider $A^T A \in \mathbb{R}^{n \times n}$. Since it is symmetric and positive semi-definite, it has orthonormal eigenvectors $v_1, \dots, v_n$ with eigenvalues $\lambda_1 \geq \dots \geq \lambda_n \geq 0$. Define $\sigma_i = \sqrt{\lambda_i}$. For $\sigma_i > 0$, let $u_i = A v_i / \sigma_i$. Orthonormality follows: $u_i^T u_j = \frac{v_i^T A^T A v_j}{\sigma_i \sigma_j} = \frac{\lambda_j \delta_{ij}}{\sigma_i \sigma_j} = \delta_{ij}$. Complete $\{u_i\}$ to an orthonormal basis of $\mathbb{R}^m$. Then $A V = U \Sigma \implies U^T A V = \Sigma$.

    \item \textbf{Low-Rank Approximation (Eckart-Young):}
    \begin{itemize}
        \item For $B_k = \sum_{i=1}^k \sigma_i u_i v_i^T$, $\| A - B_k \|_2 = \sigma_{k+1}$.
        \item For any $B$ with $\text{rank}(B) = k$, the null space $\mathcal{N}(B)$ has dimension $n-k$. The subspace $V_{k+1} = \text{span}\{v_1, \dots, v_{k+1}\}$ has dimension $k+1$.
        \item There exists a unit vector $z \in \mathcal{N}(B) \cap V_{k+1}$. Then $\|(A-B)z\|_2 = \|Az\|_2$.
        \item Writing $z = \sum_{i=1}^{k+1} \alpha_i v_i$, $\|Az\|_2^2 = \sum_{i=1}^{k+1} \sigma_i^2 \alpha_i^2 \geq \sigma_{k+1}^2$. Thus $\|A-B\|_2 \geq \sigma_{k+1}$.
    \end{itemize}
\end{enumerate}

\textbf{Intuition:}
The Singular Value Decomposition (SVD) identifies the principal axes of the transformation $A$. The 2-norm of a matrix is determined by its largest singular value. When approximating with a lower rank, the "best" strategy is to remove the components corresponding to the smallest singular values, as they contribute the least to the "energy" or "variance" of the mapping.
\end{solution}

\subsection{Team}

\begin{exercise}
\label{applied:2010:7}
Let $X_1, \dots, X_n$ be i.i.d. with distribution function $F$. Let $Y_1 < \dots < Y_n$ be the order statistics. Prove that $Z_j = F(Y_j)$ has a Beta distribution $Beta(j, n-j+1)$.
\end{exercise}

\begin{solution}
Assume $F$ is continuous, as in the source problem.  Put
\[
U_i=F(X_i).
\]
Then the probability integral transform gives $U_i\sim \operatorname{Unif}(0,1)$, and monotonicity of $F$ implies
\[
Z_j=F(Y_j)=U_{(j)},
\]
the $j$-th order statistic of $U_1,\ldots,U_n$.

For $0<z<1$, the distribution function of $U_{(j)}$ is
\[
\mathbb{P}(U_{(j)}\leq z)
=\sum_{k=j}^n {n\choose k}z^k(1-z)^{n-k},
\]
because at least $j$ of the $n$ uniform variables must fall in $[0,z]$. Differentiating gives
\[
f_{U_{(j)}}(z)
=\frac{n!}{(j-1)!(n-j)!}z^{j-1}(1-z)^{n-j},
\qquad 0<z<1.
\]
This is exactly the density of a $\operatorname{Beta}(j,n-j+1)$ random variable.
\end{solution}

\begin{exercise}
\label{applied:2010:8}
Investigate the convergence to a Normal distribution for the sequence of random variables $Y_{n,i} = \frac{3}{4b_n}(1 - X_i^2/b_n^2)I(|X_i| \le b_n)$, where $X_i$ are i.i.d. and $b_n$ is a sequence of constants.
\end{exercise}

\begin{solution}
Assume the common distribution of the $X_i$ has a density $f$ that is continuous at $0$, and consider the kernel average
\[
\hat f_n(0)=\frac{1}{n}\sum_{i=1}^n Y_{n,i}.
\]
Write
\[
K(u)=\frac{3}{4}(1-u^2)I(|u|\leq1),
\qquad
Y_{n,i}=\frac{1}{b_n}K(X_i/b_n).
\]
Then
\[
E Y_{n,i}=\int_{-1}^{1}K(u)f(b_nu)\,du\to f(0)
\]
as $b_n\to0$, because $\int_{-1}^{1}K(u)\,du=1$.

The second moment is
\[
E Y_{n,i}^2=\frac{1}{b_n}\int_{-1}^{1}K(u)^2f(b_nu)\,du
\sim \frac{f(0)}{b_n}\int_{-1}^{1}K(u)^2\,du.
\]
Since
\[
\int_{-1}^{1}K(u)^2\,du
=\frac{9}{16}\int_{-1}^{1}(1-2u^2+u^4)\,du
=\frac{3}{5},
\]
we have
\[
\operatorname{Var}(Y_{n,i})\sim \frac{3f(0)}{5b_n}.
\]

If
\[
b_n\to0,\qquad nb_n\to\infty,
\]
then the triangular-array Lindeberg condition holds because the kernel is bounded by $3/(4b_n)$, while the standard deviation of $\sum_i Y_{n,i}$ is of order $\sqrt{n/b_n}$. Hence
\[
\sqrt{nb_n}\left(\hat f_n(0)-E\hat f_n(0)\right)
\xrightarrow{d}N\left(0,\frac{3}{5}f(0)\right).
\]
If $f$ is twice continuously differentiable at $0$, then
\[
E\hat f_n(0)-f(0)=\frac{b_n^2}{10}f''(0)+o(b_n^2),
\]
so the centered CLT with $f(0)$ in place of $E\hat f_n(0)$ also holds when $nb_n^5\to0$.
\end{solution}

\begin{exercise}
\label{applied:2010:9}
Prove the minimax error bound: $\min \sup_{|\theta| \le \tau} E (\sum a_i X_i + a_{n+1} - \theta)^2 = \frac{\tau^2 n^{-1}}{\tau^2 + n^{-1}}$ for a suitable model (e.g. $X_i \sim N(\theta, 1)$).
\end{exercise}

\begin{solution}
Write the affine estimator as
\[
\delta(X)=\sum_{i=1}^n a_iX_i+a_{n+1}.
\]
Let $A=\sum_{i=1}^n a_i$. Since $X_i=\theta+\varepsilon_i$ with independent $\varepsilon_i\sim N(0,1)$,
\[
R(\theta,\delta)
=E_\theta(\delta-\theta)^2
=\sum_{i=1}^n a_i^2+\{(A-1)\theta+a_{n+1}\}^2.
\]
For fixed $A$, Cauchy's inequality gives
\[
\sum_{i=1}^n a_i^2\geq \frac{A^2}{n},
\]
with equality exactly when $a_1=\cdots=a_n=A/n$. Also, because the parameter set is symmetric, any nonzero intercept increases the larger of the two endpoint risks:
\[
\max_{\theta=\pm\tau}\{(A-1)\theta+a_{n+1}\}^2
\geq (A-1)^2\tau^2.
\]
Thus no affine estimator can have smaller worst-case risk than the best estimator of the form
\[
\delta_A=A\bar X.
\]
For this estimator,
\[
R(\theta,\delta_A)=\frac{A^2}{n}+(A-1)^2\theta^2,
\]
and the maximum over $|\theta|\leq\tau$ is
\[
\frac{A^2}{n}+(A-1)^2\tau^2.
\]
Minimizing in $A$ gives
\[
\frac{2A}{n}+2(A-1)\tau^2=0,
\qquad
A=\frac{\tau^2}{\tau^2+n^{-1}}.
\]
Substitution gives the minimax value in the affine class:
\[
\min_{a_1,\ldots,a_{n+1}}\sup_{|\theta|\leq\tau}
E_\theta(\delta-\theta)^2
=\frac{\tau^2 n^{-1}}{\tau^2+n^{-1}}.
\]
The least favorable two-point prior intuition is that the endpoints $\theta=\pm\tau$ are exactly where the optimized affine risk is largest; the symmetric shrinkage balances variance at the center against endpoint bias.
\end{solution}

\begin{exercise}
\label{applied:2010:10}
In the ``1 and 1'' foul shooting model, the shooter first takes one shot. If that shot is made, the shooter gets a second shot; if the first shot is missed, the trial ends. Let $N_0,N_1,N_2$ be the counts, in $n$ independent trials, of making $0,1,2$ baskets. Assume each shot is an independent Bernoulli trial with success probability $p$. Derive the MLE
\[
\widehat p=\frac{N_1+2N_2}{N_0+2N_1+2N_2}
\]
and find its asymptotic distribution.
\end{exercise}

\begin{solution}
Under the source model, a trial may use either one or two shots. The cell probabilities are
\[
p_0=1-p,\qquad p_1=p(1-p),\qquad p_2=p^2,
\]
and $(N_0,N_1,N_2)$ is multinomial with total $n=N_0+N_1+N_2$.  The log-likelihood, up to constants, is
\[
\ell(p)=(N_0+N_1)\log(1-p)+(N_1+2N_2)\log p.
\]
Thus
\[
\ell'(p)=\frac{N_1+2N_2}{p}-\frac{N_0+N_1}{1-p}.
\]
Setting this equal to zero gives
\[
\widehat p
=\frac{N_1+2N_2}{N_0+2N_1+2N_2}.
\]
The denominator is the random total number of shots actually attempted, not $2n$.

For the asymptotic distribution, for each trial define
\[
A=\text{number of made shots}\in\{0,1,2\},
\qquad
B=\text{number of attempted shots}\in\{1,2\}.
\]
Then
\[
\widehat p=\frac{\bar A_n}{\bar B_n}.
\]
The means are
\[
E A=p(1-p)+2p^2=p(1+p),
\qquad
E B=(1-p)+2p=1+p,
\]
so $EA/EB=p$. The delta-method influence variable for the ratio is
\[
\frac{A-pB}{1+p}.
\]
Its variance is
\[
\frac{E(A-pB)^2}{(1+p)^2}.
\]
Using the three outcomes $0,1,2$ baskets,
\[
\begin{aligned}
E(A-pB)^2
&=(1-p)p^2+p(1-p)(1-2p)^2+p^2(2-2p)^2\\
&=p(1-p^2).
\end{aligned}
\]
Therefore
\[
\sqrt n(\widehat p-p)
\xrightarrow{d}
N\left(0,\frac{p(1-p)}{1+p}\right).
\]
For finite $n$, $\widehat p$ is generally not unbiased; for example when $n=1$,
\[
E\widehat p=\frac{1}{2}p(1-p)+p^2=\frac{p(1+p)}{2}\neq p.
\]
\end{solution}

\begin{exercise}
\label{applied:2010:11}
When considering finite difference schemes approximating partial differential equations (PDEs), for example, the scheme
\[
u_j^{n+1} = u_j^n - \lambda(u_j^n - u_{j-1}^n) \tag{1}
\]
where $\lambda = \frac{\Delta t}{\Delta x}$, approximating the PDE
\[
u_t + u_x = 0, \tag{2}
\]
we are often interested in stability, namely
\[
\|u^n\| \leq C\|u^0\|, \quad n \Delta t \leq T \tag{3}
\]
for a constant $C = C(T)$ independent of the time step $\Delta t$ and the spatial mesh size $\Delta x$. Here $\|\cdot\|$ is a given norm, for example the $L^2$ norm or the $L^\infty$ norm, of the numerical solution vector $u^n = (u_1^n, u_2^n, \dots, u_N^n)$. The mesh points are $x_j = j \Delta x$, $t^n = n \Delta t$, and the numerical solution $u_j^n$ approximates the exact solution $u(x_j, t^n)$ of the PDE (2) with a periodic boundary condition.

\begin{enumerate}
    \item Prove that the scheme (1) is stable in the sense of (3) for both the $L^2$ norm and the $L^\infty$ norm under the time step restriction $\lambda \leq 1$.
    \item Since the numerical solution $u^n$ is in a finite dimensional space, Student A argues that the stability (3), once proved for a specific norm $\|\cdot\|_a$, would also automatically hold for any other norm $\|\cdot\|_b$. His argument is based on the equivalency of all norms in a finite dimensional space, namely for any two norms $\|\cdot\|_a$ and $\|\cdot\|_b$ on a finite dimensional space $W$, there exists a constant $\delta > 0$ such that
    \[
    \delta \|u\|_b \le \|u\|_a \le \frac{1}{\delta} \|u\|_b.
    \]
    Do you agree with his argument? If not, please explain the mistake made by Student A.
\end{enumerate}
\end{exercise}

\begin{solution}
\textbf{Prerequisites:}
\begin{itemize}
    \item Definition of $L^2$ and $L^\infty$ norms for vectors.
    \item Von Neumann stability analysis (Fourier analysis for finite differences).
    \item Concept of convergence and stability in the limit $\Delta x \to 0$.
\end{itemize}

\textbf{Steps:}
\begin{enumerate}
    \item \textbf{Stability Analysis:}
    \begin{itemize}
        \item \textbf{$L^\infty$ norm:} Rewrite (1) as $u_j^{n+1} = (1-\lambda) u_j^n + \lambda u_{j-1}^n$. If $0 \le \lambda \le 1$, this is a convex combination. Thus $|u_j^{n+1}| \le \max(|u_j^n|, |u_{j-1}^n|) \le \|u^n\|_\infty$. Induction gives $\|u^n\|_\infty \le \|u^0\|_\infty$.
        \item \textbf{$L^2$ norm:} Substitute $u_j^n = \hat{u}^n e^{ij\theta}$. The amplification factor is $g(\theta) = (1-\lambda) + \lambda e^{-i\theta}$. Then $|g(\theta)|^2 = 1 - 2\lambda(1-\lambda)(1-\cos\theta)$. Stability requires $|g(\theta)| \le 1$, which holds for $0 \le \lambda \le 1$.
    \end{itemize}

    \item \textbf{Student A's Mistake:} The equivalence of norms $\|u\|_a \le K \|u\|_b$ holds for a fixed dimension. However, stability requires the constant $C$ to be independent of $\Delta x$. As $\Delta x \to 0$, the dimension $N = 1/\Delta x$ goes to infinity. The equivalence constants often grow with $N$. For example, $\|u\|_2 \le \sqrt{N} \|u\|_\infty$. If a scheme is stable in $L^2$ with $C=1$, then $\|u^n\|_\infty \le \|u^n\|_2 \le \|u^0\|_2 \le \sqrt{N} \|u^0\|_\infty$. The factor $\sqrt{N}$ blows up as $\Delta x \to 0$, so $L^2$ stability does not imply $L^\infty$ stability.
\end{enumerate}

\textbf{Intuition:}
Stability is about the uniform boundedness of the numerical operator as the mesh is refined. While all norms are equivalent in any finite-dimensional space, the ratio between norms can depend on the grid size. True stability must withstand the limit $N \to \infty$.
\end{solution}

\begin{exercise}
\label{applied:2010:12}
We have the following 3 PDEs
\[
u_t + A u_x = 0, \tag{4}
\]
\[
u_t + B u_x = 0, \tag{5}
\]
\[
u_t + C u_x = 0, \quad C = A + B. \tag{6}
\]
Here $u$ is a vector of size $m$ and $A$ and $B$ are $m \times m$ real matrices. We assume $m \geq 2$ and both $A$ and $B$ are diagonalizable with only real eigenvalues. We also assume periodic initial condition for these PDEs.

\begin{enumerate}
    \item Prove that (4) and (5) are both well-posed in the $L^2$-norm. Recall that a PDE is well-posed if its solution satisfies
    \[
    \|u(\cdot, t)\| \leq C(T) \|u(\cdot, 0)\|, \quad 0 \leq t \leq T
    \]
    for a constant $C(T)$ which depends only on $T$.
    \item Is (6) guaranteed to be well-posed as well? If not, give a counter example.
    \item Suppose we have a finite difference scheme $u^{n+1} = A_h u^n$ for approximating (4) and another scheme $u^{n+1} = B_h u^n$ for approximating (5). Suppose both schemes are stable in the $L^2$-norm. If we now form the splitting scheme $u^{n+1} = B_h A_h u^n$ which is a consistent scheme for solving (6), is this scheme guaranteed to be $L^2$ stable as well? If not, give a counter example.
\end{enumerate}
\end{exercise}

\begin{solution}
\textbf{Statement-risk note.}
As stated, the first claim is true, while the two ``guaranteed'' questions in (2) and (3)
have negative answers. The missing condition in (3) is a common, grid-uniform contraction
bound for the two substeps; mere stability of the two separate one-step schemes does not
force stability of their product.

\textbf{(1) Well-posedness of (4) and (5).}
Write $A=PDP^{-1}$ with $D=\operatorname{diag}(\lambda_1,\ldots,\lambda_m)$ and
$\lambda_j\in\mathbb R$. Setting $v=P^{-1}u$ gives the decoupled system
\[
v_t+Dv_x=0,\qquad (v_j)_t+\lambda_j(v_j)_x=0.
\]
On the periodic domain each scalar component is just translated, hence
\[
\|v(\cdot,t)\|_2=\|v(\cdot,0)\|_2.
\]
Therefore
\[
\|u(\cdot,t)\|_2\leq \|P\|_2\|P^{-1}\|_2\|u(\cdot,0)\|_2,
\]
so (4) is well-posed in $L^2$. The same argument applies to $B$.

\textbf{(2) The sum need not be well-posed.}
Let
\[
A=
\begin{pmatrix}
1&1\\
0&-1
\end{pmatrix},
\qquad
B=
\begin{pmatrix}
-1&0\\
-2&1
\end{pmatrix}.
\]
Both matrices have the distinct real eigenvalues $1$ and $-1$, hence both are
diagonalizable over $\mathbb R$. But
\[
C=A+B=
\begin{pmatrix}
0&1\\
-2&0
\end{pmatrix}
\]
has characteristic polynomial $\lambda^2+2$, so its eigenvalues are
$\pm i\sqrt 2$. For the Fourier mode $e^{ikx}$, the system
$u_t+Cu_x=0$ gives
\[
\widehat u_t=-ikC\widehat u.
\]
Since $C$ has non-real eigenvalues, the matrix $-ikC$ has a real eigenvalue of size
$|k|\sqrt 2$. Thus some modes grow like $e^{|k|\sqrt 2\,t}$, and no bound depending
only on $T$ can hold uniformly in $k$. Hence (6) is not guaranteed to be well-posed.

\textbf{(3) The product of stable substeps need not be stable.}
Here ``stable'' cannot be replaced by ``each one-step matrix is power-bounded'' unless
the two bounds are compatible in one common norm. A finite-dimensional model already
shows the gap. Put
\[
S=
\begin{pmatrix}
1&0\\
0&-1
\end{pmatrix},
\qquad
T=
\begin{pmatrix}
1&1\\
0&-1
\end{pmatrix}.
\]
Then $S^2=T^2=I$, so the separate schemes $u^{n+1}=Su^n$ and $u^{n+1}=Tu^n$ are
stable. However,
\[
TS=
\begin{pmatrix}
1&-1\\
0&1
\end{pmatrix},
\qquad
(TS)^n=
\begin{pmatrix}
1&-n\\
0&1
\end{pmatrix},
\]
whose $L^2$ operator norm is unbounded as $n\to\infty$. Thus the splitting update
$u^{n+1}=TSu^n$ is not stable.

If one strengthens the hypothesis to
\[
\|A_h\|_2\leq 1,\qquad \|B_h\|_2\leq 1
\]
uniformly in $h$, then $\|B_hA_h\|_2\leq 1$ by submultiplicativity, and the product
is stable. That stronger contraction assumption is not what the problem states.
\end{solution}
